\documentclass{article}
\usepackage{graphicx} % Required for inserting images
\usepackage{biblatex} % Imports biblatex package
\addbibresource{citations.bib}
\usepackage{changepage}
\usepackage{parskip}

\title{Creating More Meaningful Explanations for Audio Deepfake Detection}
\date{2024-10-01} 
\author{Jacob LaRock\\1667321}
\linespread{1.5}

\def\Institute{Universität Siegen}
\def\KindOfWork{Bachelorarbeit}
\def\Studiengang{Wirtschaftsinformatik}
\def\Fakultaet{Fakultät III: Wirtschaftswissenschaften, Wirtschaftsinformatik und Wirtschaftsrecht}
 
\def\Title{Creating More Meaningful Explanations for Audio Deepfake Detection}
\def\Subtitle{}
 
\def\student{Jaocb LaRock}
\def\studentno{1667321}
\def\place{Siegen}
\def\Date{31.03.2025} 
\def\semester{6}
 
\def\erstpruefer{Gunnar Stevens}
\def\erstprueferMail{Gunnar.Stevens@uni-siegen.de}

\begin{document}
	\begin{titlepage}
        \begin{minipage}{0.9\linewidth}
			\centering
			\includegraphics [width=0.35\textwidth]{images/LogoSiegen}
        \end{minipage}

		\vspace{2cm}

		\centering

		{\Large\bfseries \Title\par}
		{\large\bfseries \Subtitle\par}
		\vspace{2cm}

		{\large \textbf{\KindOfWork}\par}

		\vspace{1.5cm}

		{\normalsize im Studiengang \\
			\Studiengang  im \semester. Semester \\
		}
		\vspace{0.5cm}
		{\normalsize vorgelegt von}
		\vspace{0.5cm}

		{\normalsize \textbf{\student} \\
			Matr.-Nr.: \studentno  \\}
		{am \normalsize \Date\par}

		\vspace{0.5cm}

		{\Institute} \par
		{\Fakultaet}\par

		\vspace{1cm}

    {\begin{table}[h]
			\centering
			\begin{tabular}{lll}
				{Erstprüfer:} & \erstpruefer & \erstprueferMail \\
			\end{tabular}
		\end{table}}     
		\vfill
	\end{titlepage}
    \newpage
    \section{Abstract}
    Abstract filler text blablabla!
    \section{Introduction}
    As the use of generative methods for the creation of synthetic voices becomes more widespread,
    so grows the need for reliable and usable detection methods for the protection of the security
    of people and businesses alike. In particular, the rise of deepfake technology has led to
	concerns about its potential misuse in areas such as politics, entertainment, and national
	security by, for instance, potentially allowing malicious actors to create fake audio
	recordings that appear to be genuine statements made by public figures or to create
	genuine-seeming recordings of events that never happened, which could have significant
	consequences if used to spread misinformation, from defaming individuals to creating political
	tension \cite{veerasamy_rising_2022,albahar_deepfakes_2005}. The field of Explainable AI (XAI)
	shows promise for producing useful and interpretable results from such models. Explainable AI
	refers to the ability of artificial intelligence systems, such as machine learning models and
	neural networks, to provide understandable and interpretable explanations for their decisions
	or predictions. This means that XAI systems can articulate why they made a particular
	decision, what factors influenced it, and how confident they are in their conclusion
	\cite{hind_explaining_2019}. For the detection of audio deepfakes, this could mean that the
	model can justify its classification with feature-based evidence, allowing for verification of
	the result by the end user. Much of the existing research, further discussed in the next
	section, follows the path of producing a model with the best possible benchmarks against
	datasets of samples, without taking into account if the results can be usefully interpreted,
	while existing explorations into making audio deepfake detection explainable have produced
	results that often require a lot of background knowledge to meaningfully interpret With the
	context of the research discussed in the next section, I would like to use new combinations of
	features with a novel model architecture to generate explanations for fake or real audio
	classification that are potentially more useful to their end user.
	\section{Related Work}
	\section{Experiments}
		\subsection{Dataset}
		\subsection{Black-Box Model}
			\subsubsection{Features}
			\subsubsection{Model Architecture}
		\subsection{Explanations}
	\section{Results}
		\subsection{Black-Box Model Results}
		\subsection{Explanation Results}
	\section{Discussion}
	\section{Conclusion}
	\newpage
	\printbibliography
\end{document}
